\documentclass[11pt]{article}
\usepackage[margin=1in]{geometry}
\usepackage{amsmath,amssymb}
\usepackage[T1]{fontenc}
\usepackage{hyperref}

\title{Literature Status: Sigma-Dedekind Complete Lattice Homomorphisms}
\author{Run \texttt{fa\_banach\_001}}
\date{2026-06-23}

\begin{document}
\maketitle

\paragraph{Status.}
\textbf{Literature already answered.}  The source question in
arXiv:2007.14907 asks whether the theorem that every lattice homomorphism on
an order-continuous Banach lattice attains its norm extends to
$\sigma$-Dedekind complete Banach lattices.  The later paper arXiv:2502.10165
explicitly answers this question negatively.

\paragraph{Source question.}
Dantas, Martinez-Cervantes, Rodriguez Abellan, and Rueda Zoca prove in
arXiv:2007.14907, Section 4, that every lattice homomorphism on an
order-continuous Banach lattice attains its norm.  Immediately afterwards,
around the paragraph containing their Example 4.2, they write that they do not
know whether the theorem extends to $\sigma$-Dedekind complete Banach lattices.

\paragraph{Answering paper.}
Bilokopytov, Garcia-Sanchez, de Hevia, Martinez-Cervantes, and Tradacete,
arXiv:2502.10165, Section 3, give equivalent lattice renormings of
$\ell_\infty$ and $L_\infty[0,1]$ with non-norm-attaining lattice
homomorphisms.  They explicitly state after Theorem 3.1 and its examples that
these examples answer the question from arXiv:2007.14907 in the negative,
because Dedekind completeness is an order property and is unchanged by
equivalent lattice renorming.

\paragraph{Concrete counterexample.}
Let $X=\ell_\infty$ with its usual pointwise lattice order, but give it the
equivalent lattice norm
\[
  \Vert\!\vert x\vert\!\Vert
  =
  \Vert x\Vert_\infty+\sum_{n=1}^\infty 2^{-n}|x_n|.
\]
The underlying lattice is Dedekind complete, hence $\sigma$-Dedekind complete.
Let $t\in\beta\mathbb N\setminus\mathbb N$, and let
$\delta_t\colon C(\beta\mathbb N)\to\mathbb R$ be evaluation at $t$, identified
with a free-ultrafilter limit on $\ell_\infty$.  This is a lattice
homomorphism.

For every $x\in X$,
\[
  |\delta_t(x)|\leq \Vert x\Vert_\infty\leq \Vert\!\vert x\vert\!\Vert,
\]
so $\Vert\delta_t\Vert\leq 1$ in the renormed space.  If
$u^{(m)}=\mathbf 1_{\{n>m\}}$, then $\delta_t(u^{(m)})=1$ and
\[
  \Vert\!\vert u^{(m)}\vert\!\Vert
  =
  1+\sum_{n>m}2^{-n}\longrightarrow 1,
\]
so $\Vert\delta_t\Vert=1$.  If $\delta_t$ attained its norm at some
$x$ with $\Vert\!\vert x\vert\!\Vert\leq 1$, then $|\delta_t(x)|=1$ would force
$\Vert x\Vert_\infty=1$ and
$\sum_n2^{-n}|x_n|=0$.  Hence $x=0$, contradicting $|\delta_t(x)|=1$.
Thus $\delta_t$ does not attain its norm.

\paragraph{Scope limitation.}
This answers only the $\sigma$-Dedekind-complete extension question from the
source paper.  It does not settle Conjecture 5.5 of arXiv:2007.14907 about
norm-attainment of $\widehat{x^*}$ on $FBL[E]$ versus norm-attainment of
$x^*$ on $E$.

\begin{thebibliography}{9}
\bibitem{DMCRARZ}
S.~Dantas, G.~Martinez-Cervantes, J.~D. Rodriguez Abellan, and A.~Rueda Zoca,
\emph{Norm-attaining lattice homomorphisms}, arXiv:2007.14907.

\bibitem{BGDMT}
E.~Bilokopytov, E.~Garcia-Sanchez, D.~de Hevia, G.~Martinez-Cervantes, and
P.~Tradacete,
\emph{Norm-attaining lattice homomorphisms and renormings of Banach lattices},
arXiv:2502.10165.
\end{thebibliography}

\end{document}
